\documentclass{article}

\usepackage{amsmath}
\usepackage{amssymb}
\usepackage{enumitem}
\usepackage{xcolor}

\title{Pre-Week Homework 05}
\author{Trevon Collins}
\date{\today}

\begin{document}
\maketitle

%2.1
\setcounter{section}{2}
\subsection{}
\begin{enumerate}
    \item Every real number is an even integer. $\textcolor{blue}{\text{A statement and True}}$
    \item Every even integer is a real number. $\textcolor{blue}{\text{A statement and True}}$
    \setcounter{enumi}{3}
    \item $Sets \ \mathbb{Z} \ and \ \mathbb{N} \ \textcolor{red}{\text{Not a statement}}$
\end{enumerate}

%2.5
\setcounter{subsection}{4}
\subsection{}
\begin{enumerate}
    \item $P \vee (Q \Rightarrow R)$
    \[
    \begin{array}{c|c|c|c|c}
        P & Q & R & Q \to R & P \lor (Q \to R) \\ \hline
        T & T & T &    T    &        T         \\ \hline
        T & T & F &    F    &        T         \\ \hline
        T & F & T &    T    &        T         \\ \hline
        T & F & F &    T    &        T         \\ \hline
        F & T & T &    T    &        T         \\ \hline
        F & T & F &    F    &        F         \\ \hline
        F & F & T &    T    &        T         \\ \hline
        F & F & F &    T    &        T         \\
    
    \end{array}
    \]

    \setcounter{enumi}{3}
    \item $\neg (P \lor Q) \lor (\neg P)$
    \[
    \begin{array}{c|c|c|c}
        P & Q & P \lor Q & \neg (P \lor Q) \lor (\neg P) \\ \hline
        T & T &    T     &               F               \\ \hline
        T & F &    T     &               F               \\ \hline
        F & T &    T     &               T               \\ \hline
        F & F &    F     &               T               \\
    
    \end{array}
    \]

    \setcounter{enumi}{7}
    \item $P \lor (Q \land \neg R)$
    \[
    \begin{array}{c|c|c|c|c}
        P & Q & R & Q \land \neg R & P \lor (Q \land \neg R) \\ \hline
        T & T & T &       F        &            T            \\ \hline
        T & T & F &       T        &            T            \\ \hline
        T & F & T &       F        &            T            \\ \hline
        T & F & F &       F        &            T            \\ \hline
        F & T & T &       F        &            F            \\ \hline
        F & T & F &       T        &            T            \\ \hline
        F & F & T &       F        &            F            \\ \hline
        F & F & F &       F        &            F            \\
    
    \end{array}
    \]
\end{enumerate}

%2.6
\subsection{}
\begin{enumerate}
    \setcounter{enumi}{3}
    \item $\neg (P \lor Q) = (\neg P) \land (\neg Q) $
    \[
    \begin{array}{c|c|c|c|c|c}
        P & Q & \neg P & \neg Q & \neg (P \lor Q) & (\neg P) \land (\neg Q) \\ \hline
        T & T &   F    &   F    &   \mathbf{F}    &       \mathbf{F}        \\ \hline
        T & F &   F    &   T    &   \mathbf{F}    &       \mathbf{F}        \\ \hline
        F & T &   T    &   F    &   \mathbf{F}    &       \mathbf{F}        \\ \hline
        F & F &   T    &   T    &   \mathbf{T}    &       \mathbf{T}        \\
    
    \end{array}
    \]

    \setcounter{enumi}{11}
    \item $\neg (P \Rightarrow Q)$ and $P \land \neg Q$ \\
    I think these are equivalent because the negation of $P \Rightarrow Q$
    would only be true when P is true and Q is false, which is the same
    condition for $P \land \neg Q$ to be true.
    \[
    \begin{array}{c|c|c|c|c|c}
        P & Q & \neg Q & \neg (P \Rightarrow Q) & P \land \neg Q \\ \hline
        T & T &   F    &       \mathbf{F}       &   \mathbf{F}   \\ \hline
        T & F &   T    &       \mathbf{T}       &   \mathbf{T}   \\ \hline
        F & T &   F    &       \mathbf{F}       &   \mathbf{F}   \\ \hline
        F & F &   T    &       \mathbf{F}       &   \mathbf{F}   \\
    
    \end{array}
    \]
\end{enumerate}

%2.7
\subsection{}
\begin{enumerate}
    \setcounter{enumi}{1}
    \item $\forall x \in \mathbb{R}, \exists n \in \mathbb{N}, x^{n} \geq 0 $ \\
    \textbf{For all real numbers $x$, there exists a natural number $n$ such that
    $x^{n}$ is greater than or equal to 0.} \\
    \color{blue} 
    This is true because all real numbers raised to any natural number $2n$
    will be positive and greater or equal to 0.
    \color{black}

    \setcounter{enumi}{7}
    \item $\forall \ n \in \mathbb{Z}, \exists X \subseteq \mathbb{N}, |X| = n $ \\
    \color{red} 
    This is false because the cardinality of a set cannot be negative, so 
    there is no subset of $\mathbb{N}$ with a cardinality of $-1$ for example.
    \color{black}
\end{enumerate}

%2.9
\setcounter{subsection}{8}
\subsection{}
\begin{enumerate}
    \setcounter{enumi}{4}
    \item For every $\varepsilon > 0$, there exists a $\delta > 0$ such
    that if $|x - a| < \delta \implies |f(x) - f(a)| < \varepsilon$.
    \color{blue}
    \[
    \forall \ \varepsilon \in \mathbb{N}, \exists \ \delta \in \mathbb{N},
    (|x - a|) < \delta \implies (|f(x) - f(a)|) < \varepsilon
    \]
    \color{black}


    \setcounter{enumi}{8}
    \item If $x$ is a rational number and $x \ne 0$, then tan($x$) is not a rational number. \\
    \color{blue}
    \[
    x \in \mathbb{Q} \implies \tan(x) \notin \mathbb{Q}
    \]
    \color{black}

\end{enumerate}

%2.10
\subsection{}
\begin{enumerate}
    \setcounter{enumi}{3}
    \item For every positive number $\varepsilon$, there is a positive number
            $\delta$ such that $|x - a| < \delta$ implies $|f(x) - f(a)| < \varepsilon$.
    \[ 
        \begin{aligned}
            \forall \ \varepsilon > 0, \exists \ \delta > 0&, (|x - a|) < \delta \implies (|f(x) - f(a)|) < \varepsilon \\
            \neg (\forall \ \varepsilon > 0, \exists \ \delta > 0&, (|x - a|) < \delta \implies (|f(x) - f(a)|) < \varepsilon) \\
            \exists \ \varepsilon > 0, \neg (\exists \ \delta > 0&, (|x - a|) < \delta \implies (|f(x) - f(a)|) < \varepsilon) \\
            \exists \ \varepsilon > 0, \forall \ \delta > 0&, \neg ((|x - a|) < \delta \implies (|f(x) - f(a)|) < \varepsilon) \\
            \exists \ \varepsilon > 0, \forall \ \delta > 0&, (|x - a|) < \delta \land \neg ((|f(x) - f(a)|) < \varepsilon)
        \end{aligned} 
    \]
    \fbox{
        \parbox{0.85\linewidth}{
            $
                \begin{aligned}
                \exists \ \varepsilon > 0,& \forall \ \delta > 0, (|x - a|) < \delta \land (|f(x) - f(a)|) \geq \varepsilon
                \end{aligned}
            $

            There exists a positive number $ \varepsilon$ such that 
                all positive number $\delta$ are greater than $|x - a|$
                and $|f(x) - f(a)|$ is greater than or equal to $\varepsilon$
        }
    }
    \setcounter{enumi}{7}
    \item If $x$ is a rational number and $x \ne 0$, then tan($x$) is not a rational number.
    \[
        \begin{aligned}
            \exists \ x \in \mathbb{Q}& \land x \ne 0, \tan(x) \notin \mathbb{Q} \\
            \neg (\exists \ x \in \mathbb{Q} &\land x \ne 0, \tan(x) \notin \mathbb{Q}) \\
            \forall \ x \in \mathbb{Q} &\lor x = 0, \neg(\tan(x) \notin \mathbb{Q}) \\
        \end{aligned}
    \]
    \fbox{
        \parbox{0.8\linewidth}{
            $\forall \ x \in \mathbb{Q} \lor x = 0, \tan(x) \in \mathbb{Q}$
            
            For all rational numbers $x$ or $x$ equals 0, tan($x$) is a rational number.
        }
    }
    

\end{enumerate}
\end{document}